\clearpage\thispagestyle{empty}
\pdfbookmark[0]{Sources for Moon Stars (February 1823)}{nicolai1823-sources}
\begin{center}\textbf{Sources for Moon Stars (February 1823)}\end{center}
\noindent This AI translation follows Nicolai’s \textit{Mondsterne} in \href{https://gdz.sub.uni-goettingen.de/id/PPN236007467}{\textit{Werke}, vol.~VI} (Göttingen, 1874), pp.~443--445. The paper first appeared as \textit{Berechnung der Meridiandifferenz zweier Orte, aus correspondirenden Mondsculminationen} in H.~C. Schumacher’s \textit{Astronomische Nachrichten}, vol.~II, no.~26 (February 1823), cols.~\href{https://api.digitale-sammlungen.de/iiif/image/v2/bsb10049427_00021/full/1800,/0/default.jpg}{17--18}, \href{https://api.digitale-sammlungen.de/iiif/image/v2/bsb10049427_00022/full/1800,/0/default.jpg}{19--20}, \href{https://api.digitale-sammlungen.de/iiif/image/v2/bsb10049427_00025/full/1800,/0/default.jpg}{21--22} and \href{https://api.digitale-sammlungen.de/iiif/image/v2/bsb10049427_00026/full/1800,/0/default.jpg}{23--24}. The links lead to the Bavarian State Library’s copy of the volume, published in Altona in 1824.
\par\smallskip\noindent The collected edition supplies the heading \textit{Mondsterne} and the bracketed address to Schumacher. Its row of 22 dashes marks the omission of Nicolai’s worked examples from cols.~20--24 of the first printing. These include observations on 3, 4 and 5 March and a comparison between Dorpat and Mannheim. This translation follows the collected notation, including the order of factors in the expression relating hourly motion to longitude difference, the unbarred radical, the letter $s$ in the paragraph after the probable-error formula, and the raised roman s in $0^{\mathrm{s}}1$. In col.~24 the first printing instead has $0^{\prime\prime}{,}1$. Translation and typesetting produced by AI.
\clearpage
